\subsection{Empirical stability with unbounded managed features}
\label{app:unbounded_concentration}

This subsection replaces the empirical-stability and noise proposition in
\citet{bignozzigagliardinischneider2026}. The covariance calculation uses
their alternating-block/Berbee argument, with an analytical truncation of
the feature norm. The score is controlled by a direct martingale second-moment
argument. The estimator itself is not truncated, and the remaining
operator decomposition and source-bias calculation are unchanged.

\begin{lemma}[Unbounded-feature stability and score]
\label{lem:unbounded_stability}
Assume Conditions~\ref{ec:dependence}--\ref{ec:pricing} and
$\Tr\Sigma<\infty$. Put
\[
C_{\rm sc}:=\sqrt{8K_4(1+K_4)}.
\]
For $0<\lambda\leq1$, put
\[
D_\lambda=(\Sigma+\lambda I)^{-1/2},\quad
C_\lambda=D_\lambda\Sigma D_\lambda,\quad
 d_\lambda=\Tr C_\lambda=\mathcal C_K(\lambda),
\]
\[
\Delta_{\lambda,T}=D_\lambda(\Sigma-\widehat\Sigma_T)D_\lambda,
\qquad
S_{\lambda,T}=\frac1T\sum_{t=1}^T D_\lambda m^\star_{t+1}X_t.
\]
Fix $\delta\in(0,1)$ and define
\[
\begin{split}
\ell&=\max\!\left\{1,\left\lceil
  (8C_\beta T/\delta)^{1/(\gamma+1)}\right\rceil\right\},\\
H&=\log\!\left(\frac{eT}{\delta\lambda}\right),\qquad
\tau^2=8B_X^2H,\qquad B_\lambda=\frac{1+\tau^2}{\lambda},\\
z&=\log\!\left(\frac{112(1+d_\lambda)}{\delta}\right).
\end{split}
\]
If $\ell\leq T/4$, there is a numerical constant $C$ such that, with
probability at least $1-\delta$,
\begin{equation}
\|\Delta_{\lambda,T}\|_{\op}
\leq C\!\left\{\sqrt{\frac{\ell B_\lambda z}{T}}
                  +\frac{\ell B_\lambda z}{T}\right\}
       +\frac{2B_X^2(8H+1)}{\lambda}e^{-8H},
\label{eq:unbounded_covariance}
\end{equation}
and
\begin{equation}
\|S_{\lambda,T}\|_{\Hh_K}^2
\leq\frac{2C_{\rm sc}}{\delta}\frac{d_\lambda}{T}.
\label{eq:unbounded_score}
\end{equation}
Consequently, for each fixed $\delta$ there is a class-uniform $C_\delta$
such that, for sufficiently large $T$, the sufficient condition
\begin{equation}
T^{\gamma/(\gamma+1)}\lambda
\geq C_\delta\left[\log\!\left(\frac{C_\delta T}{\lambda}\right)\right]^2
\label{eq:unbounded_admissibility}
\end{equation}
implies $\|\Delta_{\lambda,T}\|_{\op}\leq1/2$ and
\eqref{eq:unbounded_score} with probability at least $1-\delta$.
\end{lemma}

\begin{proof}
We suppress the economy index. Only the covariance is truncated in the
proof. Let
\[
A_t=(D_\lambda X_t)\otimes(D_\lambda X_t)
       \mathbf1_{\{\|X_t\|\leq\tau\}},\qquad
C_{\lambda,\tau}=\E A_t,\qquad H_t=A_t-C_{\lambda,\tau}.
\]
The process $(H_t)$ is centered and has no larger absolute-regularity
coefficients than $(Z_t,R_{t+1})$. Since
$0\preceq C_{\lambda,\tau}\preceq C_\lambda\preceq I$ and
$\|A_t\|_{\op}\leq\tau^2/\lambda$, we have
\begin{equation}
\|H_t\|_{\op}\leq B_\lambda,\qquad
\E H_t^2=\E A_t^2-C_{\lambda,\tau}^2
\preceq B_\lambda C_\lambda.
\label{eq:truncated_operator_variance}
\end{equation}
The second inequality is an operator inequality; no independence of
successive observations is used.

Let $n=\lfloor T/(2\ell)\rfloor$ and split the first $2n\ell$
observations into $2n$ consecutive blocks of length $\ell$.
Treat the odd and even blocks separately. For an odd-block sum $B_j^o$,
operator Cauchy--Schwarz gives
\[
\|B_j^o\|_{\op}\leq\ell B_\lambda,\qquad
\E(B_j^o)^2\preceq\ell^2B_\lambda C_\lambda.
\]
Indeed, for every $h$, $\|\sum_{t\in I}H_th\|^2
\leq\ell\sum_{t\in I}\|H_th\|^2$.
Berbee coupling replaces each family by independent blocks with identical
block laws. The total probability of a coupling failure for both families is
at most
\[
2n\beta(\ell)\leq C_\beta T\ell^{-(\gamma+1)}\leq\delta/8.
\]
For the independent copies, a variance majorant has operator norm at most
$n\ell^2B_\lambda$ and trace at most $n\ell^2B_\lambda d_\lambda$.
The trace form of the self-adjoint Bernstein inequality
\citep[proof of Theorem~3.1 and Section~3.2]{minsker2017}, with this
variance majorant, yields for each family, with failure probability at most
$\delta/8$,
\[
\left\|\frac1T\sum_{j=1}^n B_j^{o,*}\right\|_{\op}
\leq C\!\left\{\sqrt{\frac{\ell B_\lambda z}{T}}
                       +\frac{\ell B_\lambda z}{T}\right\}.
\]
The trace-to-variance-majorant ratio is at most $d_\lambda$; no bound on
the effective rank of the actual variance is inferred by monotonicity.
Apply the finite-dimensional inequality to increasing finite-rank
projections and pass to the limit in operator norm, as in that reference.
The operators here are trace class and bounded after truncation.
The fewer than $2\ell$ leftover observations contribute at most
$2\ell B_\lambda/T$, already absorbed because $z\geq1$.

It remains to compare the truncated and original covariances.
Lemma~\ref{lem:primitive_consequences} gives
\[
\mathbb P\!\left(\max_{t\leq T}\|X_t\|>\tau\right)
\leq2Te^{-8H}\leq\delta/8.
\]
On the complementary event, the truncated empirical covariance equals the
original empirical covariance. The population truncation bias satisfies
\[
\begin{aligned}
\|C_\lambda-C_{\lambda,\tau}\|_{\op}
&\leq\lambda^{-1}\E[\|X_t\|^2\mathbf1_{\{\|X_t\|>\tau\}}]\\
&\leq\frac{2}{\lambda}(\tau^2+B_X^2)e^{-\tau^2/B_X^2}
=\frac{2B_X^2(8H+1)}{\lambda}e^{-8H}.
\end{aligned}
\]
The tail-integration formula justifies the second line. This proves
\eqref{eq:unbounded_covariance} with total failure probability at most
$\delta/2$.

For the score, Lemma~\ref{lem:primitive_consequences} gives
\[
\E\|D_\lambda m^\star_{t+1}X_t\|^2
\leq C_{\rm sc}\Tr(D_\lambda\Sigma D_\lambda)
=C_{\rm sc}d_\lambda.
\]
For $s<t$, the score at $s$ is $\mathscr F_t$-measurable, whereas the
conditional mean of the score at $t$ given $\mathscr F_t$ is zero.
All cross inner products consequently have zero expectation. Hence
\[
\E\|S_{\lambda,T}\|^2
=\frac1{T^2}\sum_{t=1}^T
  \E\|D_\lambda m^\star_{t+1}X_t\|^2
\leq C_{\rm sc}\frac{d_\lambda}{T}.
\]
Markov's inequality proves \eqref{eq:unbounded_score} with failure
probability $\delta/2$. There is no blocking loss in this score bound.

Finally, $d_\lambda\leq B_X^2\log2/\lambda$ and, for fixed $\delta$,
$\ell\leq C_\delta T^{1/(\gamma+1)}$.
Thus the leading square-root term in \eqref{eq:unbounded_covariance} is
bounded by a constant times
\[
\left\{
\frac{[\log(C_\delta T/\lambda)]^2}
     {T^{\gamma/(\gamma+1)}\lambda}
\right\}^{1/2}.
\]
The linear term is bounded by the corresponding squared expression.
The truncation bias tends to zero uniformly for $0<\lambda\leq1$ as
$T\to\infty$. Increasing $C_\delta$ proves the final statement.
\end{proof}

\begin{corollary}[Finite-sample portfolio loss]
\label{cor:unbounded_risk}
Under the assumptions of Lemma~\ref{lem:unbounded_stability}, for each fixed
$\delta\in(0,1)$ and sufficiently large $T$ satisfying
\eqref{eq:unbounded_admissibility},
\[
Q_\varepsilon(\widehat W_{\lambda,T})-Q_\varepsilon(W^\star_\varepsilon)
\leq9R^2\lambda+C_\delta\frac{\mathcal C_K(\lambda)}{T}
\]
with probability at least $1-\delta$, uniformly over the class.
\end{corollary}
\begin{proof}
Use the error decomposition, source-bias bound and operator-interaction
bound of \citet{bignozzigagliardinischneider2026}, specialized to $r=1$.
For completeness, write
\[
\begin{split}
W_\lambda&=(\Sigma+\lambda I)^{-1}\Sigma W^\star,\qquad B=W_\lambda-W^\star,\\
V_1&=(\widehat\Sigma_T+\lambda I)^{-1}
          \frac1T\sum_t m^\star_{t+1}X_t,\\
V_2&=(\widehat\Sigma_T+\lambda I)^{-1}(\Sigma-\widehat\Sigma_T)B.
\end{split}
\]
Then $\widehat W-W^\star=B+V_1+V_2$.
Spectral calculus gives $\|\Sigma^{1/2}B\|^2\leq R^2\lambda$ and
$\|(\Sigma+\lambda I)^{1/2}B\|^2\leq2R^2\lambda$.
On the stability event,
\[
\|\Sigma^{1/2}V_1\|^2\leq4\|S_{\lambda,T}\|^2,
\qquad
\|\Sigma^{1/2}V_2\|^2\leq2R^2\lambda.
\]
The first inequality follows from $\|(I-\Delta_{\lambda,T})^{-1}\|\leq2$;
the second follows from
$\|(I-\Delta_{\lambda,T})^{-1}\Delta_{\lambda,T}\|\leq1$.
Finally, $\|u+v+w\|^2\leq3(\|u\|^2+\|v\|^2+\|w\|^2)$
and Lemma~\ref{lem:unbounded_stability} give the claim. In particular,
one may take the score constant $24C_{\rm sc}/\delta$.
\end{proof}
